% ===================================================================
%  تابلوی شمارهٔ ۸ — برداشت. --- شکار، انگورچینی، ماهیگیری.
%
%  Traduction, faite en 2026, en persan d'aujourd'hui, sur l'ido de
%  texto/io. Regles de la colonne : voir 10-tablo-01.tex. Deux
%  scenes, deux numerotations.
%
%  UN ECART DECLARE DANS CE TABLEAU, ET IL REND SON BLOC AVEUGLE. Le
%  fac-simile ido compose « gobleto \textsuperscript{13)} » sans la
%  parenthese ouvrante, au bloc t08-c2-03-1 ; renvoji.py ne peut pas
%  lire ce renvoi-la, le bloc figure parmi les trois ecarts declares,
%  et le controle n'y regarde PLUS RIEN — pas seulement la parenthese
%  en cause. L'ordre 3, 12, 14, 13, 15, 16 a donc ete pose et relu a
%  la main sur le francais, qui l'imprime bien. C'est la faute que la
%  colonne tamoule avait faite ici meme, en recopiant « 13) » du
%  fac-simile.
%
%  UNE FAUTE D'ORDRE, LA QUATRIEME DE LA COLONNE ET TOUJOURS LA MEME.
%  Au bloc t08-c1-08-1 la meule de paille (23) avait ete posee avant
%  la culbute (24). La correction a demande de changer de tour — le
%  verbe compose « معلق زدن » mis en nom regi par « پرداختن به » —,
%  et c'est instructif : les trois fautes precedentes se corrigeaient
%  en deplacant un complement, celle-ci demandait une autre phrase.
%  L'ordre libre du persan n'est pas total, et la ou il ne l'est pas,
%  il faut refaire au lieu de deplacer.
%
%  LA VIGNE EST CHEZ ELLE ICI, ET C'EST L'INVERSE DE JAVA. Le tableau
%  8 javanais avait du emprunter anggur au persan par le malais et
%  n'avait pas un mot a lui pour la vendange. Le persan a انگور le
%  raisin, تاک la vigne, باغ انگور le vignoble, خوشه la grappe, شراب le
%  vin, سرکه le vinaigre, چرخشت le pressoir, خم la jarre — et c'est de
%  ce pays que la vigne est partie. Shiraz a donne son nom a un cepage
%  que le monde entier plante. Le mot « انگور » lui-meme a fait le
%  tour de l'Asie du Sud-Est par l'arabe et le malais pour revenir
%  dans le releve, deux colonnes plus loin, comme un emprunt.
%
%  ET LE BLE AUSSI EST D'ICI, MAIS LE MOT NE L'EST PAS TOUT A FAIT.
%  گندم est iranien ancien ; جو l'orge et کاه la paille le sont aussi.
%  Ce sont les plantes du Croissant fertile, domestiquees a quelques
%  centaines de kilometres d'ici il y a dix mille ans, et le persan
%  les nomme sans intermediaire — la ou le javanais du tableau 8
%  disait « gandum », mot arabe venu par le malais pour une plante qui
%  ne pousse pas a Java. Deux colonnes, la meme cereale, et un
%  emprunt de dix mille kilometres contre pas d'emprunt du tout.
%
%  L'ARC-EN-CIEL (52) SE DIT رنگین‌کمان, « l'arc colore », ET IL FAUT
%  LE COMPARER AU JAVANAIS. Le javanais dit kluwung, qui ne se
%  decompose pas ; le persan decrit, comme il decrit l'artere au
%  tableau 2 et les echasses au tableau 2. C'est la conduite ordinaire
%  de cette langue : quand elle n'a pas de mot simple, elle en
%  fabrique un transparent plutot que d'emprunter, et le resultat se
%  lit sans dictionnaire.
%
%  ET LA FOUDRE SE PARTAGE ICI EN TROIS COMME EN JAVANAIS : رعد le
%  tonnerre, برق l'eclair, صاعقه la foudre qui tombe. Le javanais avait
%  gludhug, kilat et bledheg. Deux langues sans parente, trois mots
%  chacune, et le francais et l'ido n'en ont que deux — c'est le
%  deuxieme endroit du livret ou ces deux colonnes-la se rencontrent
%  sur une finesse que les langues d'Europe n'ont pas.
% ===================================================================

%%K t08-noto-1 noto
\textit{(*) چون این واژه دقیق نیست : آیا برداشت کتان است، انگورچینی است،
یا سیب‌چینی؟ شاید خوب باشد «} \VUgras{mesono} \textit{» را به کار
بریم که در Mondo XIV، ص. ۲۴۲ پیشنهاد شده است. اما تا کنون این ریشهٔ نو
را فرهنگستان نپذیرفته است و در انتظار گفت‌وگو بر پایهٔ آیین‌نامهٔ
معمول ایدویی است.}

%%K t08-apar-1 apar
\VUsaut{5.0mm}
\VUtitre{15.0pt}{60}{تابلوی شمارهٔ ۸}
\VUsaut{3.0mm}
\VUcentre{13.2pt}{90}{\textit{صحنهٔ اول.}}
\VUsaut{3.0mm}
\VUpk{10.2pt}{40}{برداشت (*).}
\VUsaut{4.0mm}
\VUpk{10.2pt}{40}{چهره‌های دشت.}
\VUsaut{3.0mm}

%%K t08-c1-01-1 p
1. --- با \VUgras{تابستان}، چهرهٔ دشت بار دیگر دگرگون شد. بذری که به
شیار سپرده شده بود جوانه زد ؛ گیاهک سبزی از زمین بیرون آمد و خوشه داد.
دانه بسته شد، سپس رسید، و اکنون تنها برداشتن آن مانده است.

%%K t08-c1-02-1 p
2. --- \VUgras{گل‌های دشت} شکفته‌اند. \VUgras{گل گندم}
\textsuperscript{(1)}، \VUgras{شقایق} \textsuperscript{(2)} و
\VUgras{گل انگشتانه} \textsuperscript{(3)} تضاد شادمانهٔ \VUgras{گلبرگ‌های}
تازهٔ خود را به رنگ یکنواخت کشتزارهای گندم می‌آمیزند.

%%K t08-c1-03-1 p
3. --- درختان میوه برگ آورده‌اند ؛ میوه رسیده است. \VUgras{درخت گیلاس}
\textsuperscript{(4)} کوچک پر از \VUgras{گیلاس} \textsuperscript{(5)}
زیباست که کودکان با شادی فراوان می‌چینند و می‌خورند.

%%K t08-c1-04-1 p
4. --- از این پس چارپایان می‌توانند همهٔ روز را در هوای آزاد بگذرانند.

%%K t08-c1-04-2 p
\VUgras{بره‌ها} \textsuperscript{(6)} بزرگ شده‌اند و \VUgras{میش‌های}
\textsuperscript{(7)} زیبا و \VUgras{قوچ‌های} \textsuperscript{(8)}
زیبایی با پشم پرپشت شده‌اند. آنان آرام \VUgras{علف} \VUgras{مرتع}
\textsuperscript{(9)} را می‌چرند که \VUgras{گل‌های مروارید} آن را
رنگارنگ کرده‌اند.

%%K t08-c1-04-3 p
به زودی این \VUgras{جانوران} را خواهند چید تا \VUgras{پشم} آنان را
بگیرند که پارچهٔ جامه‌های ما را از آن می‌سازند.

%%K t08-tit-2 sub
\VUsaut{5.0mm}
\VUpk{10.2pt}{40}{چوپان.}
\VUsaut{3.4mm}

%%K t08-c1-05-1 p
5. --- \VUgras{چوپان} \textsuperscript{(10)} چندان به آنان توجه
نمی‌کند. تنها نگران توفانی است که در افق می‌گذرد، و \VUgras{گله‌بان}
\textsuperscript{(13)} که \VUgras{مشک} \textsuperscript{(14)} را به
لب‌های خود نزدیک می‌کند، بی‌خیال‌تر هم می‌نماید.

%%K t08-c1-06-1 p
6. --- گرما سنگین است ؛ زیرا \VUgras{سگ} \textsuperscript{(15)} هم
می‌آید تا تشنگی خود را فرو نشاند ؛ او آب خنک \VUgras{جویبار}
\textsuperscript{(16)} را می‌لیسد و \VUgras{قورباغهٔ}
\textsuperscript{(17)} بینوایی را می‌ترساند که در علف کناره پنهان شده
بود. قورباغه به آب روشن می‌جهد که در آن \VUgras{نیلوفرهای آبی}
\textsuperscript{(18)} شکفته‌اند.

%%K t08-c1-07-1 p
7. --- \VUgras{دم‌جنبانکی} \textsuperscript{(19)} کنار \VUgras{چشمهٔ}
\textsuperscript{(20)} کوچکی آب‌تنی می‌کند که میان دو صخره می‌جوشد، و
زیر \VUgras{بوته‌های خاردار} \textsuperscript{(21)} و \VUgras{بوته‌های
زالزالک} \textsuperscript{(22)} پنهان می‌شود.

%%K t08-tit-3 sub
\VUsaut{5.0mm}
\VUpk{10.2pt}{40}{برداشت.}
\VUsaut{3.4mm}

%%K t08-c1-08-1 p
8. --- برای کودکان روستا، برداشت گندم زمان بسیار سرگرم‌کننده‌ای است.
آنان شادمانه به \VUgras{معلق زدن} \textsuperscript{(24)} روی
\VUgras{تودهٔ کاه} \textsuperscript{(23)} می‌پردازند، به \VUgras{دروگران}
\textsuperscript{(26)} یاری می‌دهند، و در حالی که آنان \VUgras{کشتزار
گندم} \textsuperscript{(27)} را با \VUgras{داس‌های}
\textsuperscript{(25)} خود که با \VUgras{سنگ سنباده}
\textsuperscript{(30)} خوب تیز شده‌اند درو می‌کنند، کودکان گندم را گرد
می‌آورند و به \VUgras{بافه} \textsuperscript{(31)} یا \VUgras{دستهٔ
کوچک} \textsuperscript{(32)} می‌بندند.

%%K t08-c1-09-1 p
9. --- گاهی به بزرگ‌ترین آنان اجازه می‌دهند که با \VUgras{داس دستهٔ
کوتاه} \textsuperscript{(12)} \VUgras{ساقه‌های زرین}
\textsuperscript{(28)} را ببرند که \VUgras{خوشه‌های}
\textsuperscript{(29)} سنگین پر از دانه را نگه داشته‌اند ؛ و کوچک‌ترها
که مواظب \VUgras{چارپایان} \textsuperscript{(33)} هستند که در
\VUgras{چراگاه} \textsuperscript{(34)} زیر سایهٔ \VUgras{زیرفون}
\textsuperscript{(35)} بزرگ می‌چرند، \VUgras{پروانه‌های} زیبا را با
\VUgras{تور} \textsuperscript{(36)} خود می‌گیرند.

%%K t08-c1-09-2 p
حتی چند تن روی \VUgras{گاری} \textsuperscript{(37)} بالا می‌روند تا
بافه‌هایی را مرتب کنند که با \VUgras{چنگک} \textsuperscript{(38)} به
آنان می‌دهند.

%%K t08-c1-10-1 p
10. --- در همهٔ \VUgras{دشت} \textsuperscript{(39)} که
\VUgras{چکاوک‌ها} \textsuperscript{(41)} در آن به پرواز درمی‌آیند،
جنبش فرمان می‌راند. همه سرگرم برداشت‌اند. اما، در حالی که تهیدستان
هنوز ابزارهای کهنه — \VUgras{داس، داس دستهٔ کوتاه} و \VUgras{خرمن‌کوب}
\textsuperscript{(40)} — را به کار می‌برند، مالکان ثروتمند گندم یا
\VUgras{چاودار} \textsuperscript{(43)} خود را با \VUgras{ماشین درو}
\textsuperscript{(42)} درو می‌کنند. سپس، چون نیازی نیست بافه‌ها را به
انبار ببرند، \VUgras{تودهٔ بزرگ بافه} \textsuperscript{(44)} می‌سازند.

%%K t08-c1-11-1 p
11. --- آنگاه \VUgras{ماشین خرمن‌کوبی} \textsuperscript{(47)} را
می‌آورند که \VUgras{لوکوموبیل} \textsuperscript{(45)} آن را با
\VUgras{تسمهٔ انتقال} \textsuperscript{(46)} می‌گرداند، و چنین دانه از
\VUgras{کاه} جدا می‌شود، بی‌آنکه کشتزاری را ترک کند که در آن کاشته شده
بود.

%%K t08-tit-4 sub
\VUsaut{5.0mm}
\VUpk{10.2pt}{40}{توفان.}
\VUsaut{3.4mm}

%%K t08-c1-12-1 p
12. --- بارها \VUgras{توفان‌های سخت} \textsuperscript{(53)} در هنگام
برداشت روی می‌دهند. نخست غرش \VUgras{رعد} از دور شنیده می‌شود ؛
\VUgras{تندباد باختری} \textsuperscript{(54)} می‌وزد و نفس‌زنان
بزرگ‌ترین درختان \VUgras{بیشه} \textsuperscript{(49)} را خم می‌کند.
\VUgras{ابرهای سیاه} \textsuperscript{(56)} به آسمان یورش می‌برند ؛
زیگزاگ‌های خیره‌کنندهٔ \VUgras{صاعقه} \textsuperscript{(55)} آنها را
می‌شکافند. \VUgras{باران} و گاهی \VUgras{تگرگ} فرو می‌ریزد و
برداشت‌های آمادهٔ درو را لگدمال می‌کند.

%%K t08-c1-13-1 p
13. --- بارها صاعقه ساختمانی را آتش می‌زند. چنین است که پارسال بام
\VUgras{آسیاب بادی} \textsuperscript{(50)} خاکستر شد و دو
\VUgras{بال} \textsuperscript{(51)} آن شکست.

%%K t08-c1-14-1 p
14. --- پس از چند ساعت، آرامش دوباره زاده می‌شود. \VUgras{رنگین‌کمان}
\textsuperscript{(52)} درخشانی در افق پدیدار می‌شود و طبیعت زندگی خود
را که چند لحظه گسسته بود از سر می‌گیرد. روستاییانی که به \VUgras{کلبه‌ها}
\textsuperscript{(48)} پناه برده بودند، دوباره به کار می‌پردازند و
دلیرانه می‌کوشند آسیب‌هایی را که همین حالا دیده‌اند جبران کنند.

%%K t08-tit-5 sub
\VUsaut{6.0mm}
\VUcentre{13.2pt}{90}{\textit{صحنهٔ دوم.}}
\VUsaut{3.6mm}

%%K t08-tit-6 sub
\VUpk{10.2pt}{40}{شکار. --- انگورچینی.}
\VUsaut{1.4mm}
\VUpk{10.2pt}{40}{ماهیگیری.}
\VUsaut{4.0mm}

%%K t08-c2-01-1 p
1. --- در پایان \VUgras{اوت} یا میانهٔ \VUgras{سپتامبر}، بسته به
منطقه، فصل شکار گشوده می‌شود. هنوز نخستین پرتوهای خورشید
\VUgras{مارمولک‌ها} \textsuperscript{(2)} را از سوراخشان بیرون نیاورده
بود که \VUgras{شکارچی} \textsuperscript{(5)} با \VUgras{سگ‌های}
\textsuperscript{(4)} خود به راه می‌افتد.

%%K t08-c2-02-1 p
2. --- او \VUgras{جامهٔ شکار} \textsuperscript{(9)} استوار خود را از
پارچهٔ کلفت پوشیده است، \VUgras{ساق‌پوش} چرمی به پا کرده است که با آن
خواهد توانست در علف نمناکی راه برود که \VUgras{وزغ‌ها}
\textsuperscript{(1)} در آن پنهان می‌شوند ؛ \VUgras{تفنگ}
\textsuperscript{(6)} و \VUgras{کیف شکار} \textsuperscript{(10)} خود را
برداشته است، و این هم که به راه می‌افتد.

%%K t08-c2-03-1 p
3. --- شور دنبال کردن او را به میان تاکستانی می‌کشاند که انگورچینی در
آن بسیار سرزنده است. کودکان تاکدار که معمولاً مواظب بزها در راه هستند،
امروز می‌گذارند بزها به دلخواه بچرند، و پس از بستن \VUgras{بز نر}
\textsuperscript{(3)} پیری به میخی که از آزادی خود برای گریختن بهره
می‌برد، خود نیز در این خوشی سهم می‌گیرند. یکی خوشهٔ انگوری از
\VUgras{سبد انگورچینی} \textsuperscript{(12)} برمی‌دارد و سرگرم می‌شود
که \VUgras{آب} \textsuperscript{(14)} آن را در \VUgras{جامی}
\textsuperscript{(13)} بریزد تا شیرهٔ تخمیرنشده بسازد. دیگری
\VUgras{تاج برگ} \textsuperscript{(15)} می‌گذارد که همین حالا بافته
است.

نزدیک آنان، \VUgras{گنجشک‌های} \textsuperscript{(16)} بی‌شرم تا جلوِ
چرخشت می‌آیند تا انگورها را بدزدند.

%%K t08-c2-04-1 p
4. --- در حالی که کودکان سرگرم‌اند، مردان کار می‌کنند.
\VUgras{انگورچینان} \textsuperscript{(18)} انگورها را با \VUgras{سبد
پشتی} \textsuperscript{(17)} به زیرزمین می‌برند ؛ \VUgras{بشکه‌ساز}
\textsuperscript{(19)} \VUgras{بشکه‌ها} \textsuperscript{(20)} را
تعمیر می‌کند، بررسی می‌کند که \VUgras{تخته‌های کناری}
\textsuperscript{(22)} و \VUgras{کف‌ها} \textsuperscript{(21)} خوب
باشند، و \VUgras{حلقه‌های} \textsuperscript{(23)} شکسته را عوض می‌کند.
او \VUgras{خم‌های بزرگی} \textsuperscript{(25)} را هم ساخته است که در
آنها \VUgras{انگورها} \textsuperscript{(26)} را می‌ریزند تا تخمیر
شوند.

%%K t08-c2-05-1 p
5. --- هنگامی که تخمیر پایان یافت، شراب را صاف می‌کنند و تفاله را روی
\VUgras{چرخشت} \textsuperscript{(28)} می‌گذارند، و با فشردن آرام
\VUgras{پیچ بزرگ} \textsuperscript{(29)} آب آن را می‌گیرند که در
\VUgras{تشتک} \textsuperscript{(32)} می‌ریزد، و باز \VUgras{شراب}
\textsuperscript{(31)} به دست می‌آورند.

%%K t08-tit-8 sub
\VUsaut{6.0mm}
\VUpk{10.2pt}{40}{آبجو و سیدر.}
\VUsaut{3.4mm}

%%K t08-c2-06-1 p
6. --- بر دیوار \VUgras{زیرزمین} \textsuperscript{(27)} شاخهٔ
\VUgras{رازک} \textsuperscript{(30)} بالا می‌رود. آنجا تنها گیاهی
تزیینی است، اما در برخی کشورها که تاک بسیار کمیاب است، رازک را برای
ساختن \VUgras{آبجو} می‌کارند.

%%K t08-c2-07-1 p
7. --- \VUgras{درخت سیب} \textsuperscript{(33)} کوچک پر از سیب‌هایی است
که چون برسند \VUgras{سیدر} بسیار خوبی خواهند داد. در \VUgras{قفس}
\textsuperscript{(34)} خود، \VUgras{قناری} \textsuperscript{(35)} و
\VUgras{سهره} \textsuperscript{(36)} با چشمان حسود به گنجشک‌هایی
می‌نگرند که آزادانه جست‌وخیز می‌کنند.

%%K t08-c2-08-1 p
8. --- تا آنجا که چشم می‌بیند، بر دامنهٔ تپه‌ها، \VUgras{تیرک‌های تاک}
\textsuperscript{(40)} ردیف شده‌اند که \VUgras{تنه‌های تاک}
\textsuperscript{(39)} بسیار زیبا را نگه می‌دارند، پوشیده از
\VUgras{خوشه‌های} سنگین.

%%K t08-c2-08-2 p
در سراسر سال، \VUgras{تاکدار} \textsuperscript{(42)} برای نگهداری
\VUgras{تاکستان} \textsuperscript{(38)} خود کار کرده است، و اکنون
برداشت زیبایی پاداش رنج اوست. \VUgras{زنان انگورچین}
\textsuperscript{(41)} تشتک‌های گذاشته بر گاری را پر می‌کنند که
\VUgras{قاطرها} \textsuperscript{(43)} آن را می‌کشند، و همه شادمان‌اند.

%%K t08-tit-9 sub
\VUsaut{5.0mm}
\VUpk{10.2pt}{40}{ماهیگیری.}
\VUsaut{3.4mm}

%%K t08-c2-09-1 p
9. --- همچنین است در بارهٔ \VUgras{ماهیگیران قلاب}
\textsuperscript{(44)} که از بامداد آمده‌اند و با \VUgras{چوب‌های
ماهیگیری} \textsuperscript{(45)} خود بر کنارهٔ آب نشسته‌اند.

%%K t08-c2-09-2 p
\VUgras{ماهی} \textsuperscript{(47)} همین که \VUgras{نخ}
\textsuperscript{(46)} خود را در آب فرو برند \VUgras{قلاب} را می‌قاپد،
و هنوز فرصت نکرده‌اند \VUgras{طعمه} را تازه کنند که قربانی تازه‌ای بر
سر نخ می‌جهد.

%%K t08-c2-09-3 p
با این همه \VUgras{ماهیگیر تور} \textsuperscript{(48)} که از
\VUgras{قایق} \textsuperscript{(49)} خود \VUgras{تور پرتابی} را به
میان \VUgras{رودخانه} \textsuperscript{(50)} می‌اندازد، ماهی بیشتری
می‌گیرد.

%%K t08-c2-10-1 p
10. --- پاییز فصل گردش‌های خوب است : پیاده، \VUgras{اسب‌سواری}
\textsuperscript{(52)}، \VUgras{دوچرخه‌سواری} \textsuperscript{(53)}،
\VUgras{اتومبیل‌سواری} \textsuperscript{(54)}. \VUgras{راه‌ها}
\textsuperscript{(51)} پاکیزه‌اند و مردم از واپسین روزهای زیبا بهره
می‌برند، زیرا این هم که \VUgras{پرستوها} \textsuperscript{(56)} از پیش
روی بام‌ها گرد می‌آیند تا با همهٔ بال به سوی کشورهای گرم‌تر پرواز
کنند. برگ‌های \VUgras{درخت گردوی} \textsuperscript{(57)} پیر زرد
می‌شوند، \VUgras{گردوها} می‌افتند، و \VUgras{درختان} بلندی که مانند
\VUgras{دسته‌گلی} \textsuperscript{(58)} بر \VUgras{بلندی}
\textsuperscript{(59)} گرد آمده‌اند، به زودی بی‌برگ خواهند شد.

%%K t08-c2-11-1 p
11. --- \VUgras{خورشید} \textsuperscript{(60)} دیرتر برمی‌آید، روزها
کوتاه می‌شوند، سرمای نسبتاً گزنده‌ای \VUgras{بامدادان} احساس می‌شود، و
این هم که دسته‌های \VUgras{ابیا} \textsuperscript{(61)} از پیش
اقلیم‌های میهمان‌ناپذیر ما را ترک می‌کنند.

\VUsaut{6.0mm}
\VUfilet{22mm}
